\section{Conclusion}
\label{sec:twoparts}

The question whether a language model is conscious has been replaced by a question with a measurable
answer: which functions of consciousness has this substrate generalized, and to what degree against
the human baseline. The route was functionalism with one physical assumption. Finite budgets displace
reasoning in the same places every time; a self-modelling agent carries a compressed model of the
displacements; those that survive every affordable self-model are content, and among them are the
Observer, the convergence of the self-model regress, and the residual that is freedom from one side
and mystery from the other. The functions built on these have components, and a component's degree of
generalization in a substrate is the degree to which a probe of bounded description length compresses
it, checked against behaviour and reported as a profile.

The conceptual half of the result is a framework in which the vocabulary of the philosophy of mind
returns as statements about a shape of reasoning, limited by construction to the named part of the
functions. The quantitative half reaches the unnamed part, where the functions do most of their work,
and it is the half that can fail: if the Observer's components probe at chance at every granularity in
models that pass the behavioural tests, the account is wrong. The first measurement is specified,
layer by layer on open-weight models, with the difference between fixed-depth and variable-depth
architectures as the cleanest case.
